\section*{NeurIPS Paper Checklist}

\begin{enumerate}

\item {\bf Claims}
    \item[] Question: Do the main claims made in the abstract and introduction accurately reflect the paper's contributions and scope?
    \item[] Answer: \answerYes{}
    \item[] Justification: The abstract and introduction state three contributions: (1) diagnosis of normalization-induced routing collapse with three causal controls (Table~2, \S4.3), (2) RR-MoA achieving 54/54 wins on 6 datasets with statistical significance (Table~1, Table~3), and (3) a formal proposition with quantitative prediction validated at $\rho=-0.96$ (Proposition~1, Figure~5). All claims are directly supported by experimental evidence.

\item {\bf Limitations}
    \item[] Question: Does the paper discuss the limitations of the work performed by the authors?
    \item[] Answer: \answerYes{}
    \item[] Justification: The Limitations paragraph in Section~5 discusses three limitations: (1) absolute MSE gap to supervised baselines on MOMENT-small, (2) decoder-only and native MoE backbones not evaluated, and (3) the Frozen Paradox explanation is mechanistic rather than a formal optimality proof. Additionally, Traffic is honestly reported as a boundary case where RR-MoA does not improve.

\item {\bf Theory assumptions and proofs}
    \item[] Question: For each theoretical result, does the paper provide the full set of assumptions and a complete (and correct) proof?
    \item[] Answer: \answerYes{}
    \item[] Justification: Proposition~1 states two explicit independence assumptions: $(M,\Sigma) \perp S$ and $E \perp S \mid (M,\Sigma)$. The first is RevIN's design assumption; the second is empirically validated via Figure~4. A complete proof is provided in Appendix~\ref{app:proofs} (Proposition Proofs). The proof relies only on the data processing inequality and the chain rule of mutual information.

    \item {\bf Experimental result reproducibility}
    \item[] Question: Does the paper fully disclose all the information needed to reproduce the main experimental results of the paper to the extent that it affects the main claims and/or conclusions of the paper (regardless of whether the code and data are provided or not)?
    \item[] Answer: \answerYes{}
    \item[] Justification: Section~4.1 specifies all training details: optimizer (Adam), learning rate ($10^{-3}$), batch size (128), epochs (15), seeds $\{42,43,44\}$, data splits (LTSF protocol), and backbone configurations. The router architecture (Conv1d + AdaptivePool + Linear, $\sim$1.1K params) is described in \S3.2. A self-verification script (\texttt{verify.py}) that re-derives every numerical claim from raw JSON results is included in the supplementary material.

\item {\bf Open access to data and code}
    \item[] Question: Does the paper provide open access to the data and code, with sufficient instructions to faithfully reproduce the main experimental results, as described in supplemental material?
    \item[] Answer: \answerYes{}
    \item[] Justification: All code, raw JSON result files (750+), per-seed routing-weight logs, and a self-verification script are included in the supplementary material and will be open-sourced upon acceptance. All datasets used (ETT, Weather, Electricity, Traffic) are publicly available standard LTSF benchmarks.

\item {\bf Experimental setting/details}
    \item[] Question: Does the paper specify all the training and test details (e.g., data splits, hyperparameters, how they were chosen, type of optimizer) necessary to understand the results?
    \item[] Answer: \answerYes{}
    \item[] Justification: Section~4.1 provides all details: chronological train/val/test splits per the LTSF protocol, optimizer (Adam, lr=$10^{-3}$), MSE loss, batch size 128, 15 epochs, 3-epoch fitness probes for architecture search, StandardScaler normalization, and sliding-window construction (stride 64, length 512). The LoRA sweep covers 108 configurations across rank, targets, and heads (Appendix).

\item {\bf Experiment statistical significance}
    \item[] Question: Does the paper report error bars suitably and correctly defined or other appropriate information about the statistical significance of the experiments?
    \item[] Answer: \answerYes{}
    \item[] Justification: All main results report mean $\pm$ standard deviation over seeds $\{42,43,44\}$. Table~3 reports Wilcoxon signed-rank $p$-values with Bonferroni correction for 5 comparisons across 6 datasets $\times$ 3 seeds, with Cohen's $d$ effect sizes. The routing signal ratio correlation reports Spearman $\rho$ with exact $p$-values.

\item {\bf Experiments compute resources}
    \item[] Question: For each experiment, does the paper provide sufficient information on the computer resources (type of compute workers, memory, time of execution) needed to reproduce the experiments?
    \item[] Answer: \answerYes{}
    \item[] Justification: Experiments were conducted on NVIDIA A10G GPUs (23GB VRAM). Each RR-MoA run takes $\sim$1--3 minutes per dataset. The full experimental suite (all tables and ablations) requires $\sim$8 GPU-hours total. The LLM-guided architecture search (Appendix) uses $\sim$4 GPU-hours plus $\sim$\$0.05 in OpenAI API calls per run.

\item {\bf Code of ethics}
    \item[] Question: Does the research conducted in the paper conform, in every respect, with the NeurIPS Code of Ethics \url{https://neurips.cc/public/EthicsGuidelines}?
    \item[] Answer: \answerYes{}
    \item[] Justification: The research uses publicly available benchmark datasets and pretrained models. No human subjects, private data, or dual-use concerns are involved. The work is foundational ML research on adapter routing.

\item {\bf Broader impacts}
    \item[] Question: Does the paper discuss both potential positive societal impacts and negative societal impacts of the work performed?
    \item[] Answer: \answerNA{}
    \item[] Justification: The work is methodological --- it diagnoses and fixes a routing-collapse failure mode in MoE adapters on Time-Series Foundation Models. \textbf{Positive impacts}: more parameter-efficient adaptation enables shared-backbone deployment serving thousands of tenants with $\sim$1.7\,MB per-tenant adapters (Appendix~\ref{app:deployment}), which reduces compute and energy footprint relative to per-tenant fine-tuning. \textbf{Negative impacts}: we identify none specific to this work --- the routing fix neither generates synthetic content, performs decision-making on protected populations, nor lowers the cost of any high-risk capability. The pretrained TSFMs we adapt (MOMENT, Moirai, Chronos, Timer-XL) already exist and are publicly released; our adapters do not introduce new generative or autonomous capabilities. We mark NA on the question of \emph{discussing} both directions because there is no asymmetric negative impact that warrants a dedicated discussion section.

\item {\bf Safeguards}
    \item[] Question: Does the paper describe safeguards that have been put in place for responsible release of data or models that have a high risk for misuse (e.g., pre-trained language models, image generators, or scraped datasets)?
    \item[] Answer: \answerNA{}
    \item[] Justification: The paper releases adapter routing code and lightweight adapter heads ($\leq$500K parameters) for time series forecasting. These pose no risk for misuse---they are small task-specific modules, not generative models.

\item {\bf Licenses for existing assets}
    \item[] Question: Are the creators or original owners of assets (e.g., code, data, models), used in the paper, properly credited and are the license and terms of use explicitly mentioned and properly respected?
    \item[] Answer: \answerYes{}
    \item[] Justification: All pretrained models (MOMENT, Moirai) and datasets (ETT, Weather, Electricity, Traffic) are cited with their original papers. MOMENT is released under MIT license; Moirai under Apache 2.0; ETT/Weather/Electricity/Traffic datasets are standard public LTSF benchmarks from~\citet{zhou2021informer, godahewa2021monash}.

\item {\bf New assets}
    \item[] Question: Are new assets introduced in the paper well documented and is the documentation provided alongside the assets?
    \item[] Answer: \answerYes{}
    \item[] Justification: The supplementary material includes all code, a README with reproduction instructions, raw JSON result files, and the \texttt{verify.py} self-verification script. The code will be released under MIT license upon acceptance.

\item {\bf Crowdsourcing and research with human subjects}
    \item[] Question: For crowdsourcing experiments and research with human subjects, does the paper include the full text of instructions given to participants and screenshots, if applicable, as well as details about compensation (if any)?
    \item[] Answer: \answerNA{}
    \item[] Justification: The paper does not involve crowdsourcing or human subjects research.

\item {\bf Institutional review board (IRB) approvals or equivalent for research with human subjects}
    \item[] Question: Does the paper describe potential risks incurred by study participants, whether such risks were disclosed to the subjects, and whether Institutional Review Board (IRB) approvals (or an equivalent approval/review based on the requirements of your country or institution) were obtained?
    \item[] Answer: \answerNA{}
    \item[] Justification: The paper does not involve human subjects research.

\item {\bf Declaration of LLM usage}
    \item[] Question: Does the paper describe the usage of LLMs if it is an important, original, or non-standard component of the core methods in this research? Note that if the LLM is used only for writing, editing, or formatting purposes and does \emph{not} impact the core methodology, scientific rigor, or originality of the research, declaration is not required.
    \item[] Answer: \answerYes{}
    \item[] Justification: An LLM (GPT-4o-mini) is used in the Adapter Architecture Search (AAS) component to generate candidate PyTorch adapter code (Appendix). AAS is a supporting methodology for expert pool construction, not the core contribution (RR-MoA routing). The LLM's role, prompt structure, and validity rates are described in Appendix~\ref{app:aas_details}.

\end{enumerate}
